% Chapter 14: Enterprise Deployment
\chapter{Enterprise Deployment}
\label{ch:enterprise}

\begin{itshape}
Your VP asks: ``Is Claude secure enough for our regulated environment?
What certifications does it have? What will it cost at 200 engineers?''
This chapter is designed to be shared with decision-makers.
\end{itshape}

\section{Security and Compliance}
\label{sec:security}

\marginnote[Enterprise]{Certifications verified Feb 2026.\sourceurl{https://trust.anthropic.com/}}

\official{Anthropic holds industry-leading certifications:}

\begin{description}
  \item[SOC 2 Type II] Verified controls for security, availability, and confidentiality.
  \item[ISO 27001] Information security management system.
  \item[ISO 42001] AI management system.
  \item[FedRAMP High] US government cloud security baseline.
  \item[HIPAA BAA] Healthcare data protection agreements available.
\end{description}

\subsection{Identity and Access Management}

\official{Enterprise-grade IAM:}
\begin{itemize}
  \item SSO via SAML 2.0 and OIDC
  \item SCIM provisioning for automated user lifecycle
  \item Role-based access control (RBAC)
  \item Audit logs for all interactions
  \item Compliance API and Analytics API for programmatic access
\end{itemize}

\subsection{Data Protection}

\official{Key data protection features:}
\begin{itemize}
  \item \textbf{Encryption}: TLS 1.2+ in transit, AES-256 at rest.
  \item \textbf{Zero Data Retention (ZDR)}: Available for API and Enterprise plans.
  \item \textbf{No training on enterprise data}: Enterprise and API data is not
    used for model training by default. Explicit opt-in required.
\end{itemize}

\marginenterprise{Enterprise/API data is never used for training by default.
This is a contractual guarantee.\sourceurl{https://privacy.claude.com/en/articles/7996885-how-do-you-use-personal-data-in-model-training}}

% -------------------------------------------------------------------
\section{Deployment Scenario: 50-Person Rollout}
\label{sec:deployment-scenario}

\practitioner{A realistic 4-week enterprise adoption timeline:}

\begin{fullwidth}
\begin{tabular}{@{}llp{8cm}@{}}
\toprule
\textbf{Week} & \textbf{Phase} & \textbf{Activities} \\
\midrule
1 & Pilot (5 developers) & Install Claude Code. Create shared CLAUDE.md.
  Establish deny rules. Identify 3 high-value workflows. \\
2 & Expand (15 developers) & Deploy managed-settings.json. Add PreToolUse hooks.
  Create team commands. Measure baseline productivity. \\
3 & Scale (50 developers) & SSO integration. Shared MCP configurations.
  Training sessions using this handbook. Monitor usage. \\
4 & Optimize & Analyze usage patterns. Tune model selection for cost.
  Implement caching strategy. Set spending caps. \\
\bottomrule
\end{tabular}
\end{fullwidth}

% -------------------------------------------------------------------
\section{Cost Optimization at Scale}
\label{sec:cost-scale}

\margincost{Prompt caching (90\%) + Batch API (50\%) = 70\%+ combined savings.\sourceurl{https://platform.claude.com/docs/en/build-with-claude/prompt-caching}}

\begin{enumerate}
  \item \textbf{Optimize CLAUDE.md for caching.}
    Stable tool definitions and system prompts cache automatically.
    Restructure to put stable content first.
  \item \textbf{Use Batch API for bulk operations.}
    Code review, documentation generation, test scaffolding---all benefit
    from 50\% batch discounts.
  \item \textbf{Right-size models.}
    \practitioner{Sonnet handles the majority of coding tasks.}
    Opus for architecture and complex reasoning.
    Haiku for formatting and simple lookups.
  \item \textbf{Set spending caps.}
    Admin controls prevent runaway costs from long-running agent sessions.
\end{enumerate}

\begin{whybox}[Caching Beats Model Downgrading]
The instinct to save money by using a cheaper model
sacrifices quality across every interaction.
Optimizing for cache hits preserves quality while reducing costs
by a larger margin. A stable CLAUDE.md that caches well
saves more than switching from Opus to Sonnet.
\end{whybox}

% -------------------------------------------------------------------
\section{Enterprise Deployments at Scale}
\label{sec:enterprise-deployments}

\begin{fullwidth}
\begin{tabular}{@{}lp{3cm}p{6cm}@{}}
\toprule
\textbf{Organization} & \textbf{Scale} & \textbf{Use Case} \\
\midrule
Deloitte\sourceurl{https://www.anthropic.com/news/deloitte-anthropic-partnership} &
  470K employees & Role-customized Claude ``personas'' for consulting \\
Accenture\sourceurl{https://www.anthropic.com/news/anthropic-accenture-partnership} &
  30K trained & CIO-focused ROI measurement \\
Snowflake\sourceurl{https://www.anthropic.com/news/snowflake-anthropic-expanded-partnership} &
  12.6K customers & Text-to-SQL integration \\
\bottomrule
\end{tabular}
\end{fullwidth}

% -------------------------------------------------------------------
\section{Competitive Positioning}
\label{sec:competitive}

\begin{decisionbox}[Claude vs.\ GPT by Task Profile]
\textbf{Choose Claude when}: long-document analysis, agentic coding workflows,
enterprise compliance requirements (FedRAMP High, ISO 42001),
API-first development.

\textbf{Choose GPT when}: consumer-facing features, multimodal applications
requiring image generation, existing GitHub Copilot integration.

\textbf{Reality}: most serious enterprise deployments use both,
choosing by task profile. The question is not ``which is better''
but ``which is better for this specific use case.''
\end{decisionbox}

% -------------------------------------------------------------------
\section{Government and Regulated Industries}
\label{sec:government}

\practitioner{Government deployment options include:}

\begin{description}
  \item[FedRAMP High] Available via Palantir FedStart
    integration.\sourceurl{https://claude.com/solutions/government}
  \item[AWS GovCloud] Cloud deployment for regulated
    requirements.\sourceurl{https://claude.com/solutions/government}
  \item[Google Vertex AI] Alternative cloud deployment for
    regulated workloads.\sourceurl{https://claude.com/solutions/government}
\end{description}

\marginnote[Warning]{Government program terms change frequently.
Verify current offering details directly at
\url{https://claude.com/solutions/government} before
making procurement decisions.}

% -------------------------------------------------------------------
\section{When to Deploy: Decision Matrix}

\begin{fullwidth}
\begin{tabular}{@{}p{3cm}p{3cm}p{3cm}p{3cm}@{}}
\toprule
\textbf{Need} & \textbf{Claude Code} & \textbf{Claude Chat} & \textbf{API / SDK} \\
\midrule
Interactive dev & Primary tool & Quick questions & --- \\
Automation & Headless mode (\cref{ch:automation}) & --- & Primary tool \\
Team standards & Settings + hooks (\cref{ch:team}) & --- & --- \\
Compliance & Enterprise plan & Enterprise plan & API + ZDR \\
Custom apps & --- & --- & Agent SDK \\
Bulk operations & Fan-out & --- & Batch API \\
\bottomrule
\end{tabular}
\end{fullwidth}

\begin{keyinsight}
Enterprise deployment is not about buying seats---it is about
building organizational capability. The technology is the easy part.
The hard part is establishing shared conventions (CLAUDE.md standards),
training teams on effective patterns (\cref{ch:team}),
and measuring ROI (usage analytics + quality metrics).
See \cref{app:maturity} for a framework to measure organizational adoption.
\end{keyinsight}

% -------------------------------------------------------------------
\begin{trybox}
Estimate your team's monthly Claude cost:
\begin{enumerate}
  \item Count developers $\times$ estimated sessions per day
    $\times$ average tokens per session.
  \item Identify your top caching opportunity (what content
    is repeated in every session?).
  \item Calculate: what would a 60\% cache hit rate save monthly?
  \item Present the estimate to your manager with the deployment
    scenario timeline from this chapter.
\end{enumerate}
\end{trybox}
